\chapter*{Sommario}
Il presente lavoro di tesi illustra le attività di ricerca e sviluppo condotte durante il periodo di tirocinio presso \textbf{\textit{InfoCube}}, focalizzandosi sull'estrazione, strutturazione e interrogazione semantica avanzata di documenti legali e normativi. Il progetto si inserisce nel contesto più ampio di \textbf{\textit{LUDOVICA}} (\textit{Language Understanding for Document Optimization and Validation In Civil Administration}), un motore generativo basato sull'Intelligenza Artificiale volto a supportare la pubblica amministrazione locale nell'elaborazione e stesura di atti amministrativi.

Il problema tecnico principale affrontato in questo lavoro è la trasformazione di documenti PDF complessi in una base di conoscenza solida e interrogabile dalle macchine (\textbf{\textit{Knowledge Graph}}). A tale scopo, la rigida gerarchia tipica degli atti normativi (Titoli, Capi, Articoli, Commi) è stata modellata e salvata all'interno del database a grafo \textbf{\textit{Neo4j}}\cite{neo4j}, permettendo di preservare fedelmente le relazioni strutturali e concettuali dei testi originali.

Per abilitare il recupero intelligente e contestuale delle informazioni, è stato implementato un sistema di \textbf{\textit{Semantic Search}}. Questa fase ha incluso il confronto di molteplici \textbf{\textit{modelli di embedding}}, con l'obiettivo di osservarne il comportamento sul linguaggio giuridico italiano. Le scelte architetturali sono state valutate mediante un benchmark diagnostico di domande annotate, confrontando un approccio di ricerca vettoriale classica, basato sulla \textit{cosine similarity}, con una strategia che impiega le relazioni giuridiche del grafo per riordinare i risultati.

Sul benchmark diagnostico di 70 domande, la strategia arricchita dalla struttura aumenta la completezza soprattutto nei casi in cui la risposta coinvolge disposizioni collegate, mentre a cutoff fisso introduce un costo nei casi puramente semantici. Il confronto con un encoder contestuale mostra inoltre che il \textit{late chunking} non rende ridondante il contributo delle relazioni esplicite nelle condizioni valutate. I risultati descrivono pertanto un beneficio condizionato alla tipologia della domanda e non costituiscono una stima delle prestazioni su traffico reale.
